\begin{figure}[p]
    \centering
    \begin{minipage}[t]{0.48\textwidth}
        \centering
        \vspace{0pt}
        \textbf{(a) GravNet baseline}\par\vspace{3mm}
        \begin{tikzpicture}[
            >=Latex,
            node distance=5mm,
            every node/.style={font=\normalsize},
            stage/.style={draw, rounded corners, align=center,
                         minimum height=14mm, text width=62mm,
                         fill=RoyalBlue!8},
            block/.style={draw, rounded corners, align=center,
                         minimum height=28mm, text width=62mm,
                         fill=ForestGreen!11},
            arrow/.style={->, thick}
        ]
            \node[stage] (g-input) {
                \textbf{Mixed detector nodes}\\
                ECal and trigger-pad rows,\\
                \(\mathbf{X}_{\mathrm{cat}}
                \in\mathbb{R}^{N_{\Sigma}\times 8}\)
            };
            \node[stage, below=of g-input] (g-projection) {
                \textbf{Input projection}\\
                \(\operatorname{Linear}(8,128)\)\\
                \(\rightarrow\operatorname{ReLU}\)
            };
            \node[block, below=of g-projection] (g-gravnet) {
                \textbf{GravNet residual block}\\
                \(\boldsymbol{\times 4}\)\\
                \(\operatorname{GravNetConv}(128,128;\)\\
                \(d_s=4,\ d_p=128,\ k=16)\)\\
                \(\rightarrow\operatorname{ReLU}
                \rightarrow\operatorname{Dropout}(0.1)\)\\
                \(\rightarrow\text{residual addition}\)\\
                \(\rightarrow\operatorname{BatchNorm1d}(128)\)
            };
            \node[stage, below=of g-gravnet] (g-head) {
                \textbf{Per-node classification head}\\
                \(\operatorname{Linear}(128,C)\), with\\
                \(C=2\) for \(2e\) and \(C=3\) for \(3e\)
            };
            \node[stage, below=of g-head] (g-output) {
                \textbf{ECal-supervised output}\\
                Select ECal-node logits;\\
                trigger-pad nodes provide context only
            };

            \draw[arrow] (g-input) -- (g-projection);
            \draw[arrow] (g-projection) -- (g-gravnet);
            \draw[arrow] (g-gravnet) -- (g-head);
            \draw[arrow] (g-head) -- (g-output);
        \end{tikzpicture}
    \end{minipage}%
    \hfill
    \begin{minipage}[t]{0.48\textwidth}
        \centering
        \vspace{0pt}
        \textbf{(b) Transformer baseline}\par\vspace{3mm}
        \begin{tikzpicture}[
            >=Latex,
            node distance=5mm,
            every node/.style={font=\normalsize},
            stage/.style={draw, rounded corners, align=center,
                         minimum height=14mm, text width=62mm,
                         fill=RoyalBlue!8},
            block/.style={draw, rounded corners, align=center,
                         minimum height=38mm, text width=62mm,
                         fill=Orange!14},
            arrow/.style={->, thick}
        ]
            \node[stage] (t-input) {
                \textbf{Padded mixed-detector tokens}\\
                ECal and trigger-pad rows,\\
                \(\mathbf{X}_{\mathrm{pad}}\in
                \mathbb{R}^{B\times N_{\max}\times 8}\)
            };
            \node[stage, below=of t-input] (t-projection) {
                \textbf{Input projection}\\
                \(\operatorname{Linear}(8,128)\)\\
                \(\rightarrow\operatorname{ReLU}\)\\
                \(\rightarrow\operatorname{Linear}(128,128)\)
            };
            \node[block, below=of t-projection] (t-encoder) {
                \textbf{Pre-layer-normalised}\\
                \textbf{Transformer encoder block}\\
                \(\boldsymbol{\times 3}\)\\
                \(\operatorname{LayerNorm}(128)\)\\
                \(\rightarrow\text{four-head full self-attention}\)\\
                \(\rightarrow\operatorname{Dropout}(0.1)\)\\
                \(\rightarrow\text{residual addition}\)\\
                \(\operatorname{LayerNorm}(128)\)\\
                \(\rightarrow\text{FFN }128\!\rightarrow\!256
                \!\rightarrow\!128\)\\
                (ReLU, dropout \(0.1\))\\
                \(\rightarrow\text{residual addition}\)
            };
            \node[stage, below=of t-encoder] (t-finalnorm) {
                \textbf{Final encoder normalisation}\\
                \(\operatorname{LayerNorm}(128)\)
            };
            \node[stage, below=of t-finalnorm] (t-head) {
                \textbf{Per-token classification head}\\
                \(\operatorname{Linear}(128,C)\), with\\
                \(C=2\) for \(2e\) and \(C=3\) for \(3e\)
            };
            \node[stage, below=of t-head] (t-output) {
                \textbf{ECal-supervised output}\\
                Select ECal-token logits;\\
                trigger-pad tokens provide context only
            };

            \draw[arrow] (t-input) -- (t-projection);
            \draw[arrow] (t-projection) -- (t-encoder);
            \draw[arrow] (t-encoder) -- (t-finalnorm);
            \draw[arrow] (t-finalnorm) -- (t-head);
            \draw[arrow] (t-head) -- (t-output);
        \end{tikzpicture}
    \end{minipage}
    \caption{Architectures of the selected (a) GravNet and (b)
    pre-layer-normalised Transformer baselines. For GravNet, the learned-space
    dimension is \(d_s=4\), the propagated-feature dimension is \(d_p=128\),
    and each learned neighbourhood contains \(k=16\) nodes. The Transformer
    has model width 128, feed-forward width 256, and four attention heads of
    dimension 32.}
    \label{fig:baseline_architectures}
    \label{fig:gravnet_baseline_architecture}
    \label{fig:transformer_baseline_architecture}
\end{figure}
